%TODO:
% - Modificare procedimento accettazione pull request (accetta solo il verificatore, non più il responsabile)
        % Per la repo della documentazione possiamo scrivere che può pushare chiunque e accetta solo il verificatore, mentre per la repo del codice pushano solo i programmatori e accetta solo il verificatore. 
% - Aggiungere che associamo le azioni dei verbali con le issue di GitHub tramite la colonna "ID Issue Github" nella tabella delle azioni
% - Aggiungere che il nome del file è il nome del documento + versione
% - Dire che le versioni iniziano da V.0.1

\documentclass[10pt, a4paper]{article}

% Parametri che modificano il file main.tex
% Le uniche parti da cambiare su main.tex sono:
% - vari \vspace tra sezioni
% - tabella azioni da intraprendere
% - sezione altro

\def\data{2023-11-09}
\def\oraInizio{17:00}
\def\oraFine{18:30}
\def\luogo{Sede di AzzurroDigitale: \\ Via della Croce Rossa, 42 (PD)}

\def\tipoVerb{Esterno} % Interno - Esterno

\def\nomeResp{Feltrin E.} % Cognome N.
\def\nomeVer{Campese M.} % Cognome N.
\def\nomeSegr{Ciriolo I.} % Cognome N.

\def\nomeAzienda{Azzurro Digitale}
\def\firmaAzienda{azzurrodigitale.png}
\def\firmaResp{emanuele.png} % nome Responsabile

\def\listaPartInt{
Bresolin G.,
Campese M.,
Ciriolo I.,
Dugo A.,
Feltrin E.,
Michelon R.,
Orlandi G.
}

\def\listaPartEst{
\textit{Azzurro Digitale:},
Caliendo G.,
D'Avanzo C.,
Picello M., \\
\textit{Coding Cowboys:},
Cecchin A.,
Dolzan M.,
Ferraioli F.,
Giacomuzzo F.,
Lago L.,
Menon G.,
Nordio A.
}

\def\listaRevisioneAzioni {x}

\def\listaOrdineGiorno {
{Visita della sede dell'azienda AzzurroDigitale;},
{Breve presentazione dell'azienda da parte di Picello C., specialista in People\&Culture, con introduzione ai valori aziendali di AzzurroDigitale;},
{Possibilità di porre alcune domande generiche sul progetto ai referenti Caliendo G. e D'Avanzo C., con seguente approfondimento dei vari temi toccati.\\ All'incontro ed alle discussioni hanno partecipato anche i membri del gruppo \textit{Coding Cowboys}.}
}

\def\listaDiscussioneInterna {
{Discussione 1;},
{Discussione 2;},
}

\newcommand{\domris}[2]{\textbf{#1}\\#2}

\def\listaDiscussioneEsterna {
\domris
{Quali riferimenti dei membri del gruppo è preferibile inoltrare all'azienda?
}
{Sono necessari gli account che devono avere l'accesso al \textit{repository} \textit{GitHub} fornito dall'azienda, un indirizzo e-mail del team per l'accesso al canale \textit{Slack} fornito dall'azienda ed i singoli indirizzi e-mail dei componenti;
},
\domris
{Quando sarà possibile discutere nello specifico dell'analisi dei requisiti del progetto?
}
{L'azienda aveva già in programma un incontro in data 2023-11-15 della durata di 1.30h per la discussione dei requisiti del progetto con ciascun gruppo preso singolarmente. I gruppi hanno garantito entrambi disponibilità a tali incontri;
},
\domris
{L'azienda in precedenza ha consigliato l'adozione di un framework di tipo \textit{agile} (nello specifico di \textit{Scrum}) per lo sviluppo del progetto: dal momento che i membri del team non hanno mai avuto la possibilità di sperimentarlo, l'azienda propone qualche consiglio per la sua applicazione?
}
{L'azienda ha fornito una spiegazione dei vari ruoli all'interno di un framework \textit{Scrum} generico fornendo al team le basi per iniziare ad apprenderlo ed applicarlo, specificando che si tratta una metodologia complessa che si apprende con maggior successo durante la sua applicazione. Viene sottolineata l'importanza di avere un solido sistema di \textit{backlog} per l'accertamento dei progressi del progetto;
},
\domris
{Quali criteri utilizzerà l'azienda per la valutazione del successo del progetto?
}
{È stato spiegato che anche questo aspetto del progetto sarà analizzato in futuro con il team stesso, poiché al momento non vi è ancora un'idea precisa;
},
\domris
{Cosa intende di preciso l'azienda riferendosi ad un sistema di \textit{Bug Reporting} (e alla sua relativa documentazione)?
}
{L'azienda ha chiarito che la documentazione non è l'interesse principale e che come sistema di \textit{Bug Reporting} è sufficiente servirsi di un\textit{ Issue Tracking System}.
}
} 




\def\listaDecisioni {x}

\usepackage{style}
\usepackage{headerfooter}

\title{\titolo}
\author{SWEetCode}

\begin{document}

% PRIMA PAGINA
\begin{titlepage}
    \thispagestyle{empty}
    \begin{tikzpicture}[remember picture, overlay]
        % TRIANGOLI
        \draw[fill=secondarycolor, secondarycolor] (current page.north west) -- (current page.south west) -- (8.8, -28);
        \draw[fill=primarycolor, primarycolor] (-3, 5) -- (4, -13.6) -- (11, 5);

        % LOGO
        \node [xshift=-5cm, yshift=25cm] (logo) at (current page.south east) {\includegraphics[width=6.5cm]{img/logo.png}};

        % SWEETCODE - DATE
        \node [anchor=north east, align=right, xshift=-1.2cm, yshift=20.5cm, text=black] (sweetcode) at (current page.south east) {\fontsize{32pt}{36pt}\selectfont SWEetCode};
        \draw[line width=4pt, lightcol] ([xshift=-3cm, yshift=-0.37cm]sweetcode.south west) -- ([yshift=-0.37cm]sweetcode.south east);
        \node [anchor=north east, align=right, xshift=-1.2cm, yshift=18.7cm, text=black] (date) at (current page.south east){\fontsize{24pt}{24pt} \selectfont Verbale \tipoVerb};

        % NOME FILE
        \node [anchor=north east, text width=15cm, align=right, xshift=-1.2cm, yshift=17cm, text=black] (titolo) at (current page.south east){\fontsize{48pt}{48pt}\textbf{\data}};

        % BOX DATI PARTECIPANTI
        \ifthenelse{\equal{\tipoVerb}{Esterno}}{
            \node[anchor=north east, xshift=-1.2cm, yshift=14.5cm, minimum width=8cm] (box) at (current page.south east){};
        }{
            \node[anchor=north east, xshift=-1.2cm, yshift=12.5cm, minimum width=8cm] (box) at (current page.south east){};
        }

        % RESPONSABILE
        \node[anchor=north west, align=left] (dati2) at (box.north west) {\fontsize{15pt}{15pt}\selectfont \textbf{Responsabile}};
        \draw[line width=4pt, lightcol] (dati2.south west) -- ([xshift=8cm]dati2.south west);
        \node[anchor=north west, align=left] (dati21) at (dati2.south west){\fontsize{13pt}{13pt}\selectfont \nomeResp};

        % VERIFICATORE
        \node[anchor=north west, yshift=-1cm, align=left] (dati3) at (dati21.north west) {\fontsize{15pt}{15pt}\selectfont \textbf{Verificatore}};
        \draw[line width=4pt, lightcol] (dati3.south west) -- ([xshift=8cm]dati3.south west);
        \node[anchor=north west, align=left] (dati31) at (dati3.south west){\fontsize{13pt}{13pt}\selectfont \nomeVer};

        % SEGRETARIO DI RIUNIONE
        \node[anchor=north west, yshift=-1cm, align=left] (dati4) at (dati31.north west) {\fontsize{15pt}{15pt}\selectfont \textbf{Segretario di Riunione}};
        \draw[line width=4pt, lightcol] (dati4.south west) -- ([xshift=8cm]dati4.south west);
        \node[anchor=north west, align=left] (dati41) at (dati4.south west){\fontsize{13pt}{13pt}\selectfont \nomeSegr};
        
        % UNIPD - SWE
        \node [xshift=4.4cm, yshift=2.3cm, draw, secondarycolor, text=white] (uni) at (current page.south west) {\fontsize{20pt}{20pt} \selectfont Università di Padova};
        \node [xshift=0.65cm, yshift=0.7cm, draw, secondarycolor, text=white, below=of uni] (corso) {\fontsize{20pt}{20pt}\selectfont Ingegneria del Software};

        % FIRMA AZIENDA
        \ifthenelse{\equal{\tipoVerb}{Esterno}}{
            \draw[line width=4pt, lightcol] ([xshift=-1.2cm, yshift=7cm]current page.south east) -- ([xshift=-8cm, yshift=7cm]current page.south east);
            \node [xshift=-4.8cm, yshift=8cm] (logo) at (current page.south east) {\includegraphics[width=6cm]{img/firme/\firmaAzienda}};
            \node[anchor=north west, xshift=12.9cm, yshift=6.65cm, align=left] at (current page.south west)
            {\fontsize{13pt}{13pt}\selectfont L'azienda: \nomeAzienda};
        }

        % FIRMA CODING COWBOYS
        \draw[line width=4pt, lightcol] ([xshift=-1.2cm, yshift=4.35cm]current page.south east) -- ([xshift=-8cm, yshift=4.35cm]current page.south east);
        \node [xshift=-4.8cm, yshift=5.5cm] (logo) at (current page.south east) {\includegraphics[width=6cm]{img/firme/menon.png}};
        \node[anchor=north west, xshift=12.9cm, yshift=4cm, align=left] at (current page.south west)
        {\fontsize{13pt}{13pt}\selectfont Coding Cowboys: Menon G.};

        % FIRMA
        \draw[line width=4pt, lightcol] ([xshift=-1.2cm, yshift=1.8cm]current page.south east) -- ([xshift=-8cm, yshift=1.8cm]current page.south east);
        \node [xshift=-4.8cm, yshift=2.45cm] (logo) at (current page.south east) {\includegraphics[width=6cm]{img/firme/\firmaResp}};
        \node[anchor=north west, xshift=12.9cm, yshift=1.45cm, align=left] at (current page.south west)
        {\fontsize{13pt}{13pt}\selectfont Il Responsabile: \nomeResp};
        
    \end{tikzpicture}
\end{titlepage}

% REGISTRO DELLE VERSIONI
{\renewcommand{\arraystretch}{1.5}
\section*{Registro delle versioni}
\begin{tabularx}{\textwidth}{c|c|c|c|X}
\textbf{Versione} & \textbf{Data} & \quantities{\textbf{Responsabile di}\\\textbf{stesura}}& \textbf{Revisore} & \quantities{\textbf{Dettaglio e}\\\textbf{motivazioni}} \\
\hline
v1.0.0(3) & $2023-11-17$ & Michelon R. & Bresolin G. & Correzione paragrafo scopo del documento.\\
\hline
v0.0.1(22) & $2023-11-13$ & Bresolin G. & Campese M. & Aggiornamento struttura di documento: modifica ordine registro delle versioni dalla più recente e aggiunta voce ”motivazioni”; verbali: introduzione \#issue nella tabella delle azioni.\\
\hline
v0.0.1(21) & $2023-11-13$ & Bresolin G. & Campese M. & Aggiornamento documentazione: struttura di documento: presenza obbligatoria registro delle versioni e rimozione versione;  norme tipografiche: aggiornamento versioni. Aggiunta configuration managment: version control.\\
\hline
v0.0.1(20) & $2023-10-30$ & Bresolin G. & Michelon R. & Processi di supporto: configuration management: \textit{ITS} e pull request; Processi organizzativi: \textit{GitHub Teams}.\\ 
\hline
v0.0.1(19) & $2023-10-30$ & Campese M. & Michelon R. & Documentazione: strutture di documento: posta elettronica.\\
\hline
v0.0.1(18) & $2023-10-30$ & Dugo A. & Bresolin G. & Processi organizzativi: comunicazione.\\
\hline
v0.0.1(13) & $2023-10-28$ & Ciriolo I. & Campese M. & Strumenti.\\
\hline
v0.0.1(8) & $2023-10-27$ & Bresolin G. & Ciriolo I. & Processo di supporto: documentazione e verifica.\\
\hline
v0.0.1(5) & $2023-10-26$ & Campese M. & Orlandi G. &  Introduzione.\\

\end{tabularx}}

\newpage

% INDICE
\tableofcontents
\newpage

% INTRODUZIONE
\section{Introduzione}
\subsection{Scopo del documento}
\paragraph{}Questo documento si prefigge lo scopo di normare e documentare un corretto \textit{Way of Working} che verrà seguito lungo l'intero svolgimento del progetto.\\
	Tutti i membri del team sono tenuti a visionare tale documento e a seguire le norme contenute in esso. Tali norme permettono di perseguire e migliorare la coerenza e l'omogeneità del lavoro e dei file prodotti.
	Viene inoltre fornito uno storico delle modifiche e delle versioni del documento per seguire un approccio incrementale nella sua redazione, facilitandone i cambiamenti nel corso del tempo.
%\subsection{Scopo del prodotto}
%\subsection{Glossario}
%\paragraph{}Per evitare ambiguità e incomprensioni relative al linguaggio e ai termini utilizzati nella documentazione relativa al progetto viene presentato un Glossario. I termini ambigui o specifici presenti nello stesso, verranno identificati con un pedice |g|. (DA VALUTARE COME RICONOSCERE IL TERMINE )

%\subsection{Riferimenti}
%    \subsubsection{Riferimenti normativi}
%    \subsubsection{Riferimenti informativi}

% PROCESSI PRIMARI
\newpage
\section{Processi primari}

% PROCESSI DI SUPPORTO
\newpage
\section{Processi di supporto}
Nella seguente sezione vengono presentati i processi di supporto allo sviluppo del progetto.

\subsection{Configuration management}
In questa sezione vengono presentate tutte le attività che coinvolgono la gestione delle configurazioni del software realizzate da SWEetCode.

\subsubsection{Issue tracking system (ITS)} % TODO: Scrivere che se le issue sono relative ad un'azione presente in un verbale, nella tabella azioni è riportato l'id della issue
SWEetCode utilizza l'\textit{Issue Tracking System} fornito da \textit{GitHub}.\\
In questo sistema di tracciamento ogni issue viene dotata di:
\begin{itemize}
    \item Titolo: un titolo breve ma esplicativo dell'attività associata alla issue;
    \item Id numerico univoco: tale id viene generato in automatico dall'ITS;
    \item Assegnatari: i membri del team che a cui è affidata la responsabilità di risolvere la issue;
    \item Descrizione: una breve descrizione inerente alla issue associata;
    \item Data: la data di crezione della issue;
    \item Stato: gli stati possibili di una issue sono \textit{Todo}, \textit{In Progress} ed infine \textit{Done};
    \item Labels: etichette che permettono di organizzare e categorizzare le issue presenti nel sistema, rendendo più semplice e immediato il tipo di task a loro associato;
    \item Data: data di creazione della issue;
    \item Milestone: ogni issue può essere associata ad una determinata milestone del repository.
\end{itemize}
Ogni issue creata dal gruppo viene inserita all'interno di una \textit{Board} che permette di tracciare in modo rapido lo stato interno di ogni issue.

\paragraph{Automatizzazione chiusura issue}Una automatizzazione utilizzata dal team consiste nella chiusura automatica delle issue: ogni volta che un membro del team effettua un commit nel repository finalizzato alla chiusura di una issue, il contenuto di tale commit dovrà prevedere all'interno del suo contenuto l'istruzione "\textit{Close \#id\_issue}".
Le conseguenze di tale operazione a seguito dell'approvazione della pull request da parte del responsabile di progetto e successiva operazione di merge nel main branch sono:
\begin{itemize}
    \item modifica dello stato della issue al valore \textit{Done};
    \item spostamento della issue all'interno della \textit{Board} sotto la sua voce "\textit{Done}".
\end{itemize}

\subsubsection{Pull requests} %TODO: Cambiare chi accetta pull request da Responsabili a Verificatori
A seguito dell'adozione dei \textit{GitHub Teams}, SWEetCode propone l'utilizzo delle pull requests per permettere un rilascio controllato: l'unico \textit{Team} dotato dei permessi necessari all'approvazione di una pull request sono i \textit{Responsabili}.\\
Una volta proposta la pull request il responsabile attuale del progetto, membro del \textit{Team Responsabili}, dovrà effettuare un ulteriore verifica sulle modifiche proposte e, a seguito delle sue valutazioni, dovrà decidere se declinare o accettare la pull request, con conseguente merge nel branch all'interno del quale la pull request è stata sollevata in caso di esito positivo.

\subsubsection{Version control}
Il sistema di \textit{Version control} adottato dal team prevede una versione nel seguente formato: vX.Y.Z(build).
\begin{itemize}
    \item v: sta per versione, rimane immutato lungo tutte le versioni;
    \item X: numero che indica la versione principale di riferimento. Viene incrementato quando viene superata una fase di revisione di avanzamento;
    \item Y: numero che viene incrementato con l'introduzione di nuove features o di miglioramenti significativi;
    \item Z:  numero che indica piccoli cambiamenti rispetto alla versione vX.Y, come il fix di bug o il caricamento di nuovi documenti all'interno del repository (ad esclusione di verbali, sia interni che esterni);
    \item (build): numero che indica le build di progetto eseguite nella versione vX.Y.Z. Tale valore viene incrementato a seguito di modifiche/aggiunte nella documentazione.
\end{itemize}
Il sistema di numerazione delle versioni ha come prima versione il valore "v0.0.1(0)".\\
Ogni volta che viene incrementato il valore X, i valori Y, Z e (build) ripartono dai loro valori iniziali, ovvero "0".\\
Ogni volta che viene incrementato il valore Y, il valore Z e (build) ripartono dal valore ”0”.\\
Ogni volta che viene incrementato il valore Z, il valore (build) riparte dal valore "0".

\subsection{Documentazione}
Tutte le attività di sviluppo sono accompagnate da una loro documentazione, con l’obiettivo di permettere una consultazione semplice e rapida alle informazioni inerenti al prodotto e alle attività stesse, svolgendo inoltre un ruolo di storicizzazione e di supporto alla manutenzione.
    \subsubsection{Ciclo di vita}
    Ogni documento prodotto da SWEetCode presenta le seguenti fasi:
        \begin{enumerate}
            \item Assegnazione: $\\$
            Ogni documento viene assegnato ad uno o più responsabili di stesura affiancati a loro volta da uno o più revisori.\\
            Una volta presa la decisione per i precedenti ruoli viene aperta su \textit{GitHub} una issue a loro associata relativa alla  realizzazione del documento, la quale viene immediatamente inserita all'interno della Board del progetto. Ogni issue viene inoltre etichettata da una label relativa al tipo di documento da produrre;
            \item Stesura: $\\$
            La stesura viene realizzata dal segretario di riunione (nel caso in cui il documento in questione sia un verbale interno o esterno) o dai responsabili di stesura.
            La produzione della documentazione in formato \textit{Latex} viene gestita tramite la piattaforma \textit{Overleaf}, la quale rende disponibile gli strumenti adatti ad una eventuale stesura sincronizzata e collaborativa tra i membri del team;
            \item Verifica e validazione: $\\$
            La fase di verifica viene realizzata dal revisore del documento: in questo stadio il documento viene sottoposto ad un controllo di contenuto e sintassi, tenendo conto delle linee guida definite dal team nelle "Norme di progetto". 
            Se il documento è di tipo esterno, una volta eseguita la verifica da parte del revisore, l'atto viene condiviso con l'ente terzo in causa, per essere sottoposto ad una sua ulteriore verifica e validazione.
            In caso di esito positivo viene eseguito un commit nel repository \href{https://github.com/sweetcode-team/Documentation}{\textcolor{black}{\textit{Documentation}}}; su \textit{GitHub} che porta alla creazione di una \textit{pull request} (il commit a lui associato dovrà chiudere la issue aperta relativa al documento in oggetto attraverso l'automazione \textit{Close \#id\_issue} ).
            In caso di esito negativo si ritorna invece alla fase precedente in modo da produrre una documentazione che tiene conto delle precisazioni presentate dal revisore e/o dall'eventuale parte esterna;
            \item Pubblicazione: $\\$
            La pubblicazione costituisce l'ultima fase del ciclo di vita del documento e avviene solamente nel caso in cui la fase di verifica abbia avuto un riscontro positivo.
            In questa fase il responsabile attuale di progetto fornisce un'ulteriore conferma alla documentazione prodotta accettando la \textit{pull request} precedentemente aperta all'interno del repository \href{https://github.com/sweetcode-team/Documentation}{\textcolor{black}{\textit{Documentation}}}; del team su \textit{GitHub}.
        \end{enumerate}
    \subsubsection{Template}
    Il team ha deciso di produrre un template grafico ricorrente all'interno della documentazione.\\
    Ogni documento presenta la medesima pagina di copertina, con una grafica realizzata tramite \textit{Adobe Illustrator} ed in seguito disposta nella pagina tramite \textit{Overleaf}, dove viene esposto il nostro logo, il nostro nome del team e i membri del team.\\
    Le seguenti pagine di ciascun documento presentano un:
        \begin{itemize}
            \item header: intestazione contenente il nostro logo, il nome del documento ed il nostro nome;
            \item footer: piè di pagina contenente i link alla nostra email e al nostro account \href{https://sweetcode-team.github.io/}{\textcolor{black}{\textit{GitHub}}}; ed infine la numerazione di pagina.
        \end{itemize}
    \subsubsection{Struttura di documento}
    Tutti i documenti prodotti da SWEetCode presentano una struttura comune:
    \begin{itemize}
        \item Pagina di copertina: descritta nella sezione Template precedente;
        \item Registro delle versioni: presente in tutti i documenti ad eccezione dei verbali. Questo registro utilizzato per tenere traccia delle varie versioni permette di comprendere in modo rapido chi ha realizzato determinate sezioni della documentazione, quando sono state realizzate, cosa è stato aggiunto e perché (sotto la voce "Dettaglio e motivazioni"). Il registro presenta le versioni ordinate a partire dalla versione più recente; 
        \item Indice: presente nel caso in cui si abbia un documento di notevole dimensione, dotato di sezioni. Il suo obiettivo è quello di facilitare e agevolare l'accesso ad un determinato contenuto trattato nel documento;
        \item Contenuto: il contenuto vero e proprio del documento.
    \end{itemize}
    %\textbf{Verbali}$\\$
    \paragraph{Verbali}Fanno eccezione a questa struttura i verbali, sia esterni che interni, che oltre a presentare nella prima pagina di copertina i ruoli inerenti alla sua produzione al posto dell'elenco dei membri del team e una sezione apposita per la firma del responsabile di riunione (e in caso di verbale esterno anche la firma da parte di un rappresentate dell'azienda), offre una struttura particolare.\\
    Il contenuto presenta infatti le seguenti sezioni:
    \begin{itemize}
        \item Intestazione: sezione contenente le informazioni inerenti a data, ora e luogo della riunione;
        \item Partecipanti: sezione contenente l'elenco in ordine alfabetico dei partecipanti interni ed eventualmente esterni alla riunione e una voce apposita per gli assenti;
        \item Revisione delle azioni: sezione contenente un elenco cronologico delle azioni assegnate nelle riunioni precedenti, con informazioni in merito al progresso, problematiche riscontrate ed eventuale conclusione di tale azione;
        \item Ordine del giorno: elenco cronologico e approssimativo delle tematiche trattate all'interno della riunione;
        \item Discussione: sezione contente l'elenco cronologico delle tematiche affrontate, trattate con una descrizione accurata delle informazioni emerse;
        \item Decisioni: decisioni prossime e/o future da intraprendere o intraprese a seguito della riunione affrontata;
        \item Azioni da intraprendere: tabella contente le azioni assegnate durante l'incontro e da svolgere nel futuro immediato. Ogni voce è accompagnata dal numero della issue associata (se presente) nell'ITS su \textit{GitHub}, un incaricato, da un revisore che dovrà verificare il suo operato e da una data di scadenza entro la quale l'azione da svolgere dovrà essere ritenuta completata;
        \item Altro: sezione contente informazioni aggiuntive quali riferimenti esterni e la data della prossima riunione (nel caso in cui essa sia stata stabilita).
    \end{itemize}    
    % \textbf{Posta elettronica}$\\$
    \paragraph{Posta elettronica}La posta elettronica rappresenta uno strumento fondamentale di comunicazione asincrona con proponenti e committente.\\ 
    Per uniformare lo stile delle mail, vengono descritti di seguito delle norme a cui attenersi nella loro stesura:
    \begin{itemize}
        \item Oggetto: campo da riempire obbligatoriamente in maniera breve e concisa;
        \item Firma: ogni mail viene conclusa con la firma digitale del gruppo collocata automaticamente;
        \item Forma: le informazioni e le richieste più importanti vanno posizionate all'inizio del messaggio. La lunghezza del testo della mail non deve essere eccessivo, onde evitare la perdita di visibilità di informazioni preziose; qualora fosse necessario un testo complesso e denso di informazioni, ci si deve servire dell'uso di paragrafi ed elenchi per strutturarlo;
        \item Stile: è buona norma seguire lo stile del mittente. La mail sarà dotata di formule di saluto iniziale e finale quali ad esempio: "\textit{Gentile..}", "\textit{Spettabile..}", "\textit{Cordiali saluti}" e "\textit{Distinti saluti}".
    \end{itemize}
    
    \subsubsection{Norme tipografiche}
        % \textbf{Citazioni tecnologiche e strumentali}$\\$
        \paragraph{Citazioni tecnologiche e strumentali}Ogni tecnologia e/o strumento menzionato all'interno della documentazione deve essere scritta in \textit{corsivo}.
        
        % \textbf{Date}$\\$
        \paragraph{Date}Tutte le date presenti nella documentazione prevedono il seguente formato: AAAA-MM-GG.
        
        % \textbf{Elenchi}$\\$
        \paragraph{Elenchi}Tutti gli elenchi presenti nella documentazione prevedono un ordine:
            \begin{itemize}
                \item alfabetico: sono puntati che non presentano alcuna correlazione con l'aspetto temporale;
                \item cronologico: sono elenchi numerati e descrivono il seguirsi di avvenimenti o argomenti nel tempo.
            \end{itemize}
        In tutti gli elenchi ogni punto termina con il ";", fatta eccezione per l'ultimo elemento che termina con il ".".
        
        % \textbf{Menzioni}$\\$
        \paragraph{Menzioni}Ogni menzione di una persona, interna od esterna al team, avviene nel seguente formato: Cognome N. (dove la lettera N sta per l'iniziale del nome).
        
        % \textbf{Nome del documento}$\\$
        \paragraph{Nome del documento}I nomi dei documenti prevedono la lettera iniziale maiuscola seguiti dal resto dei caratteri in case minuscolo e dalla versione del documento stesso.

        I verbali, sia interni che esterni, fanno eccezione e presentano un titolo composta dalla data in cui sono stati svolti nel seguente formato: AAAA-MM-GG.
        
        % \textbf{Nome del gruppo} $\\$
        \paragraph{Nome del gruppo}Il nome del gruppo presenta il seguente formato: "SWEetCode".
        
        % \textbf{Stile tipografico}$\\$
        \paragraph{Stile tipografico}La formattazione del testo dei documenti segue rigorosamente un unico stile tipografico: \textit{Poppins}.
        
        % \textbf{Versioni}$\\$ 
        \paragraph{Versioni}Ogni documento viene associato nel nostro \href{https://sweetcode-team.github.io/}{sito GitHub} alla versione del progetto per cui risulta essere ancora in vigore.\\
        All'interno del documento, nella prima riga del registro delle versioni, viene invece presentata la versione del progetto in cui è avvenuta l'ultima modifica al documento.\\
        L'incremento della versione del progetto a seguito di elaborazioni nella documentazione viene gestito attraverso la seguente norma: l'aggiunta di documentazione non precedentemente presente all'interno del repository o la modifica della documentazione precedentemente presente all'interno del repository porta all'incremento del solo valore "(build)".

\subsection{Verifica}
Ogni produzione realizzata da SWEetCode viene sottoposta ad una fase di verifica che permette la sua validazione.
Onde evitare conflitti di interesse, tale compito viene assegnato ad uno o più membri differenti da coloro che hanno conseguito alla realizzazione del prodotto in questione.

\subsubsection{Verifica dei documenti}
Ogni documento prodotto dal team passa per una o più fasi di verifica prima di essere pubblicato ufficialmente nel repository \textit{GitHub}.\\
Tale fasi vengono realizzate solamente in seguito alla consegna del documento da parte dei responsabili di stesura, i quali ritengono che in quel momento il prodotto da loro realizzato sia completato.\\
La verifica dei documenti prevede che venga assicurato:
\begin{itemize}
    \item Il rispetto delle norme tipografiche stabilite nelle \textit{Norme di progetto} del team;
    \item Il rispetto della struttura definita nelle \textit{Norme di progetto} del team;
    \item Il rispetto del contenuto, il quale deve essere chiaro, coerente, preciso ed esaustivo;
    \item Il rispetto della sintassi italiana e della correttezza ortografica dei termini utilizzati.
\end{itemize}


% PROCESSI ORGANIZZATIVI
\newpage
\section{Processi organizzativi}
\subsection{Comunicazione}
SWEetCode ha deciso di mantenere una comunicazione frequente tra i membri del gruppo. In particolare, di seguito si trova la modalità di svolgimento per le comunicazioni interne e le applicazioni utilizzate.
\subsubsection{Comunicazioni interne}
% \textbf{Comunicazioni sincrone}\\
\paragraph{Comunicazioni sincrone}Il gruppo ha deciso che gli incontri si svolgeranno  in presenza oppure attraverso la piattaforma \textit{Discord}. Per gli incontri in presenza, essi verranno decisi alcuni giorni prima in modo tale da poter consentire la presenza di tutti i membri del gruppo. Per quanto riguarda \textit{Discord}, invece, si opta per una comunicazione asincrona anche pochi minuti prima del collegamento. In entrambe le situazioni, il gruppo utilizza un approccio libero alla discussione e improntato alla crescita.

% \textbf{Comunicazioni asincrone}\\
\paragraph{Comunicazioni asincrone}Il gruppo utilizzerà \textit{Whatsapp} e \textit{Notion} per le comunicazioni asincrone. In particolare è stata creata una community su \textit{Whatsapp}, composta da un gruppo principale per comunicazioni informali, una bacheca nella quale verranno inviati i messaggi più importanti dei quali è necessario prendere visione, ed altri canali in cui i componenti del gruppo si organizzano per task collaborativi.

\subsection{GitHub Teams}
SWEetCode, attraverso un profilo \textit{Organizations}, ha deciso di utilizzare una struttura interna ai membri del gruppo che usufruisce dei \textit{Teams} di \textit{GitHub}: ogni ruolo presente all'interno del progetto porta alla creazione di un \textit{Team} avente il medesimo nome.\\
I \textit{Teams} presentati sono dunque:
\begin{itemize}
    \item Amministratori;
    \item Analisti;
    \item Progettisti;
    \item Programmatori;
    \item Responsabili;
    \item Verificatori.
\end{itemize}
Lo strumento \textit{Teams} consente di gestire i vari ruoli nel migliore dei modi, permettendo di assegnare responsabilità e permessi specifici a determinate posizioni piuttosto che ai membri del gruppo. Questa mansione altrimenti risulterebbe dispendiosa e disagevole a causa della rotazione dinamica interna dei ruoli. 


% STRUMENTI
\newpage
\section{Strumenti}
Nella seguente sezione vengono presentati tutti gli strumenti utili al team per lo svolgimento di ogni attività.

\subsection{Adobe Illustrator}
\href{https://www.adobe.com}{https://www.adobe.com}\\\\
Software utilizzato per la creazione dei template ufficiali del team, per la generazione del logo e per la scelta della palette di colori da usare. 

\subsection{Canva}
\href{https://www.canva.com}{https://www.canva.com}\\\\
Sito utilizzato per la produzione delle slide per i Diari di Bordo. É stato creato un account condiviso utilizzabile da tutti i membri.

\subsection{Discord}
\href{https://discord.com}{https://discord.com}\\\\
Piattaforma utilizzata per lo svolgimento di riunioni formali ed informali del team, attraverso la creazione di un server apposito con un canale dedicato ai messaggi di testo e diverse sale dedicate alle conversazioni vocali.

\subsection{Notion}
\href{https://www.notion.so/product}{https://www.notion.so/product}\\\\
Spazio di lavoro utilizzato per l’organizzazione delle prime task divise per membro del gruppo e per tenere traccia delle attività e dei progressi portati avanti da ogni componente. Usato anche per l’assegnazione chiara dei ruoli, come calendario condiviso del gruppo e per il salvataggio link utili in modo da favorire l’accesso in maniera veloce a tutti i documenti ed i siti necessari.

\subsection{Overleaf}
\href{https://www.overleaf.com}{https://www.overleaf.com}\\\\
Sito utilizzato per la scrittura di tutta la documentazione ufficiale in linguaggio LaTeX, attraverso la creazione di un account condiviso per tutti i membri del team.


\end{document}
